\documentclass[10pt,a4paper]{amsart}
\usepackage[margin=19mm]{geometry}
\usepackage{amsmath,amssymb,mathtools}
\usepackage[scr=rsfs,cal=euler]{mathalfa}
\usepackage{xfrac}
\usepackage[unicode,hidelinks,bookmarks=false]{hyperref}
\newcommand{\ie}{i.e.\ }
\newcommand{\cf}{cf.\ }
\newcommand{\resp}{respectively\ }
\newcommand{\Subsection}{\subsection}
\providecommand{\nobreakdash}{\nobreak\hbox{-}}
\setlength{\emergencystretch}{2em}
\multlinegap0pt
\usepackage{amssymb}
\usepackage{mathtools}
\usepackage[shortlabels]{enumitem}
\newtheorem{lemma}{Lemma}
\newcommand{\calH}{\mathcal H}
\newcommand{\calL}{\mathcal L}
\newcommand{\calR}{\mathcal R}
\DeclareMathOperator{\dist}{dist}
\newcommand{\divh}{\nabla_h\!\cdot}
\newcommand{\nabh}{\nabla_h}
\newcommand{\norm}[2]{\left\|#1\right\|_{#2}}
\title{Strict 2.5D Shadows for One-Component Navier--Stokes Regularity}
\author{Runlong Yu}
\date{Source: arXiv:2606.11720v1 (CC BY 4.0)}
\begin{document}
\maketitle

\noindent\emph{Source and licence.} Strict 2.5D Shadows for One-Component Navier--Stokes Regularity. Version arXiv:2606.11720v1 (CC BY 4.0), distributed by arXiv under the Creative Commons Attribution 4.0 International licence, http://creativecommons.org/licenses/by/4.0/, as recorded in the arXiv licence metadata for this exact version and shown on the abstract page https://arxiv.org/abs/2606.11720v1. Full licence notice: \texttt{licenses/arxiv-2606.11720-CC-BY-4.0.txt}.\par\smallskip

\noindent\textit{Editorial scope.} Counted unit: the article's lemma pI:lem:localized-qual-rough-shadow (source0.tex lines 1309--1399). The article's complete original proof of it is delivered with it (source0.tex lines 1401--1678, one \texttt{proof} environment of 278 lines). No mathematics is written, repaired or re-derived: the statement and the proof are the article's own lines, in the article's own notation, and no retained line is substituted or re-flowed.  No further source-local context is needed: every cross-reference of every retained line resolves inside the two delivered spans.  Nothing was omitted from any retained span, so the package carries no omission record.  No retained line contains a citation, so the wrapper carries no bibliography and no bibliography entry.  No retained line contains a figure environment, an image-inclusion command or drawing code, so no image is delivered and the package declares no asset dependency.  Editorial formatting only, in addition: an AMS \texttt{amsart} preamble that restates the article's own declarations and macros used by the retained lines, namely the environments \texttt{lemma} on separate counters, together with the standard packages those lines need.  The retained environments print in this document as Lemma 1 in order of appearance, whereas the article prints the same items with the numbering of its own full document.  The complete native source of the article as distributed in the arXiv e-print for this exact version is bundled unchanged as \texttt{sources/202570/source0.tex}.
\begin{lemma}[Localized qualitative rough-shadow closure]
\label{pI:lem:localized-qual-rough-shadow}
Let
\[
        0<\theta<\theta_1<\rho<1,
        \qquad
        M\ge1 .
\]
Let \((U^{(n)},P^{(n)})\) be a sequence on \(Q_\rho\), where
\[
        U^{(n)}=(U_h^{(n)},0),
        \qquad
        U_h^{(n)}=(U_1^{(n)},U_2^{(n)}).
\]
Assume
\[
        \divh U_h^{(n)}=0
        \qquad\text{in }\mathcal D'(Q_\rho),
\]
and
\begin{equation}\label{pI:eq:rough-shadow-uniform-bounds}
\begin{aligned}
&\norm{U^{(n)}}{L^\infty_tL^2_x(Q_\rho)}
+
\norm{\nabla U^{(n)}}{L^2(Q_\rho)}
+
\norm{U^{(n)}}{L^3(Q_\rho)}
\nonumber\\
&\qquad
+
\norm{P^{(n)}-(P^{(n)})_{B_\rho}(t)}{L^{3/2}(Q_\rho)}
\le M .
\end{aligned}
\end{equation}
Assume that \((U^{(n)},P^{(n)})\) satisfies the approximate horizontal limiting equation
\begin{equation}\label{pI:eq:approx-horizontal-limiting}
\partial_t U_h^{(n)}
-\Delta U_h^{(n)}
+\divh(U_h^{(n)}\otimes U_h^{(n)})
+\nabh P^{(n)}
=
\divh R^{(n)}+F^{(n)}
\end{equation}
in \(\mathcal D'(Q_\rho)\), with
\[
        R^{(n)}\to0
        \qquad\text{in }L^{3/2}(Q_\rho),
\]
and
\[
        F^{(n)}\to0
        \qquad\text{in }
        L^{3/2}\bigl((-\rho^2,0);W^{-1,3/2}(B_\rho)\bigr).
\]
Finally assume the pressure-compatibility defect tends to zero:
\[
        \partial_3P^{(n)}\to0
        \qquad\text{in }
        L^{3/2}\bigl((-\rho^2,0);W^{-1,3/2}(B_\rho)\bigr).
\]
Then, after passing to a subsequence, there exists a strict limiting-system solution \((V,Q)\) in \(Q_\theta\), namely
\[
        V=(V_h,0),
        \qquad
        \divh V_h=0,
        \qquad
        \partial_3Q=0,
\]
and
\[
        \partial_tV_h-\Delta V_h
        +\divh(V_h\otimes V_h)+\nabh Q=0,
\]
such that
\[
        U^{(n)}\to V
        \qquad\text{strongly in }L^3(Q_\theta),
\]
and
\[
        \inf_{h^{(n)}\in\calH(Q_\theta)}
        \norm{P^{(n)}-Q-h^{(n)}}{L^{3/2}(Q_\theta)}
        \to0 .
\]
Consequently,
\[
        \dist_{\rm har}
        \bigl((U^{(n)},P^{(n)}),\calL^{\rm str}(Q_\theta)\bigr)
        \to0 .
\]
\end{lemma}

\begin{proof}
We first normalize the pressures.  Replacing \(P^{(n)}\) by
\[
        P^{(n)}-(P^{(n)})_{B_\rho}(t)
\]
does not change \eqref{pI:eq:approx-horizontal-limiting} and does not change the harmonic-pressure distance, because functions of time are spatially harmonic.  Thus we assume throughout the proof that
\[
        (P^{(n)})_{B_\rho}(t)=0
\]
for almost every \(t\).

The uniform bounds imply that \(U^{(n)}\) is bounded in
\[
        L^\infty_tL^2_x(Q_\rho)\cap L^2_tH^1_x(Q_\rho).
\]
From the approximate equation, \(\partial_tU_h^{(n)}\) is bounded in a negative Sobolev space, for instance in the sum
\[
        L^{3/2}\bigl((-\rho^2,0);W^{-1,3/2}(B_\rho)\bigr)
        +
        L^2\bigl((-\rho^2,0);H^{-1}(B_\rho)\bigr).
\]
Indeed, the terms
\[
        \Delta U_h^{(n)},\quad
        \divh(U_h^{(n)}\otimes U_h^{(n)}),\quad
        \nabh P^{(n)},\quad
        \divh R^{(n)},\quad
        F^{(n)}
\]
are uniformly bounded in these negative spaces.  By the Aubin--Lions compactness lemma, after passing to a subsequence,
\[
        U^{(n)}\to V
        \qquad\text{strongly in }L^2_{\rm loc}(Q_\rho).
\]
The energy bound also gives \(U^{(n)}\) bounded in \(L^{10/3}_{\rm loc}(Q_\rho)\).  Interpolating between the strong \(L^2_{\rm loc}\) convergence and this uniform \(L^{10/3}_{\rm loc}\) bound yields
\[
        U^{(n)}\to V
        \qquad\text{strongly in }L^3_{\rm loc}(Q_\rho).
\]
Since every \(U^{(n)}\) has zero third component,
\[
        V=(V_h,0).
\]
Passing to the limit in \(\divh U_h^{(n)}=0\) gives
\[
        \divh V_h=0 .
\]

The normalized pressure sequence is bounded in \(L^{3/2}(Q_\rho)\).  Hence, after passing to a further subsequence,
\[
        P^{(n)}\rightharpoonup Q
        \qquad\text{weakly in }L^{3/2}(Q_\rho).
\]
The pressure-compatibility assumption gives
\[
        \partial_3P^{(n)}\to0
        \qquad\text{in }
        L^{3/2}\bigl((-\rho^2,0);W^{-1,3/2}(B_\rho)\bigr).
\]
Passing to the weak limit yields
\[
        \partial_3Q=0
        \qquad\text{in }\mathcal D'(Q_\rho).
\]

We now pass to the approximate horizontal equation.  Since
\[
        U_h^{(n)}\to V_h
        \qquad\text{strongly in }L^3_{\rm loc}(Q_\rho),
\]
we have
\[
        U_h^{(n)}\otimes U_h^{(n)}
        \to
        V_h\otimes V_h
        \qquad\text{strongly in }L^{3/2}_{\rm loc}(Q_\rho).
\]
The residuals satisfy
\[
        R^{(n)}\to0,
        \qquad
        F^{(n)}\to0
\]
in the stated norms.  Therefore, passing to the limit in \eqref{pI:eq:approx-horizontal-limiting} gives
\[
        \partial_tV_h-\Delta V_h
        +\divh(V_h\otimes V_h)
        +\nabh Q=0
\]
in \(\mathcal D'(Q_\rho)\).  Thus \((V,Q)\) satisfies the strict limiting system.

It remains to prove strong pressure convergence modulo spatially harmonic functions.  Choose
\[
        \chi\in C_c^\infty(B_\rho),
        \qquad
        \chi\equiv1\quad\text{on }B_{\theta_1}.
\]
Define the local Calderon--Zygmund pressures
\[
        \Pi^{(n)}
        =
        \calR_a\calR_b
        \bigl(\chi U_a^{(n)}U_b^{(n)}\bigr),
        \qquad a,b\in\{1,2\},
\]
and
\[
        \Pi
        =
        \calR_a\calR_b
        \bigl(\chi V_aV_b\bigr).
\]
Since
\[
        U_h^{(n)}\otimes U_h^{(n)}
        \to
        V_h\otimes V_h
        \qquad\text{strongly in }L^{3/2}(Q_{\theta_1}),
\]
the Calderon--Zygmund estimate gives
\[
        \Pi^{(n)}\to\Pi
        \qquad\text{strongly in }L^{3/2}(Q_{\theta_1}).
\]

We claim that \(P^{(n)}-\Pi^{(n)}\) is asymptotically harmonic.  Taking horizontal divergence of \eqref{pI:eq:approx-horizontal-limiting} and using \(\divh U_h^{(n)}=0\), we obtain
\[
        \Delta_hP^{(n)}
        =
        -\partial_a\partial_b(U_a^{(n)}U_b^{(n)})
        +\partial_a\partial_bR_{ab}^{(n)}
        +\partial_aF_a^{(n)} .
\]
Hence, in \(B_{\theta_1}\), where \(\chi\equiv1\),
\[
        -\Delta P^{(n)}
        =
        \partial_a\partial_b(U_a^{(n)}U_b^{(n)})
        -\partial_a\partial_bR_{ab}^{(n)}
        -\partial_aF_a^{(n)}
        -\partial_3^2P^{(n)} .
\]
On the other hand,
\[
        -\Delta\Pi^{(n)}
        =
        \partial_a\partial_b(U_a^{(n)}U_b^{(n)})
        \qquad\text{in }B_{\theta_1}.
\]
Therefore
\[
        -\Delta(P^{(n)}-\Pi^{(n)})
        =
        -\partial_a\partial_bR_{ab}^{(n)}
        -\partial_aF_a^{(n)}
        -\partial_3^2P^{(n)} .
\]
Consequently,
\begin{align*}
&\norm{\Delta(P^{(n)}-\Pi^{(n)})}
{L^{3/2}_tW^{-2,3/2}_x(Q_{\theta_1})}
\\
&\qquad\le
C\norm{R^{(n)}}{L^{3/2}(Q_\rho)}
+
C\norm{F^{(n)}}{L^{3/2}_tW^{-1,3/2}_x(B_\rho)}
+
C\norm{\partial_3P^{(n)}}{L^{3/2}_tW^{-1,3/2}_x(B_\rho)} .
\end{align*}
The right-hand side tends to zero.  Thus
\[
        \Delta(P^{(n)}-\Pi^{(n)})\to0
        \qquad\text{in }
        L^{3/2}\bigl((-\theta_1^2,0);W^{-2,3/2}(B_{\theta_1})\bigr).
\]

We now use an elementary harmonic-projection fact.  If \(f^{(n)}\) is bounded in \(L^{3/2}(Q_{\theta_1})\) and
\[
        \Delta f^{(n)}\to0
        \qquad\text{in }
        L^{3/2}\bigl((-\theta_1^2,0);W^{-2,3/2}(B_{\theta_1})\bigr),
\]
then there exist \(h_0^{(n)}\in\calH(Q_\theta)\) such that
\[
        \norm{f^{(n)}-h_0^{(n)}}{L^{3/2}(Q_\theta)}\to0 .
\]
Indeed, let \(D^{-1}_{\theta_1}\) denote the zero-Dirichlet transposition inverse for \(-\Delta\) on \(B_{\theta_1}\).  For almost every \(t\), set
\[
        w^{(n)}(\cdot,t)
        :=D^{-1}_{\theta_1}\bigl(-\Delta f^{(n)}(\cdot,t)\bigr),
\]
that is,
\[
        -\Delta w^{(n)}(\cdot,t)
        =
        -\Delta f^{(n)}(\cdot,t)
        \quad\text{in }B_{\theta_1},
        \qquad
        w^{(n)}(\cdot,t)=0
        \quad\text{on }\partial B_{\theta_1}.
\]
The inverse is a bounded linear map
\[
        W^{-2,3/2}(B_{\theta_1})\longrightarrow L^{3/2}(B_{\theta_1}),
\]
so \(w^{(n)}\) is a strongly measurable Bochner function of time and
\[
        \norm{w^{(n)}}{L^{3/2}(B_{\theta_1})}
        \le
        C\norm{\Delta f^{(n)}}{W^{-2,3/2}(B_{\theta_1})}
\]
for a.e. \(t\).  Then \(h_0^{(n)}=f^{(n)}-w^{(n)}\) is spatially harmonic in \(B_{\theta_1}\), and restriction to \(B_\theta\) gives the desired convergence after integration in time.

Applying this fact with
\[
        f^{(n)}=P^{(n)}-\Pi^{(n)},
\]
we obtain \(h_0^{(n)}\in\calH(Q_\theta)\) such that
\[
        \norm{P^{(n)}-\Pi^{(n)}-h_0^{(n)}}{L^{3/2}(Q_\theta)}\to0 .
\]

It remains to identify the limiting harmonic correction.  Since the limit \((V,Q)\) satisfies the strict limiting system, taking horizontal divergence gives
\[
        -\Delta_hQ
        =
        \partial_a\partial_b(V_aV_b).
\]
Because \(\partial_3Q=0\), this is also
\[
        -\Delta Q
        =
        \partial_a\partial_b(V_aV_b).
\]
Inside \(B_{\theta_1}\),
\[
        -\Delta\Pi
        =
        \partial_a\partial_b(V_aV_b).
\]
Therefore
\[
        \Delta(\Pi-Q)=0
        \qquad\text{in }B_{\theta_1}
\]
for almost every \(t\).

Define
\[
        h^{(n)}=h_0^{(n)}+\Pi-Q .
\]
Then \(h^{(n)}\in\calH(Q_\theta)\).  Moreover,
\begin{align*}
P^{(n)}-Q-h^{(n)}
&=P^{(n)}-Q-h_0^{(n)}-\Pi+Q
\\
&=P^{(n)}-\Pi-h_0^{(n)}
\\
&=\bigl(P^{(n)}-\Pi^{(n)}-h_0^{(n)}\bigr)
+\bigl(\Pi^{(n)}-\Pi\bigr).
\end{align*}
Both terms tend to zero in \(L^{3/2}(Q_\theta)\).  Hence
\[
        \norm{P^{(n)}-Q-h^{(n)}}{L^{3/2}(Q_\theta)}\to0 .
\]
Together with
\[
        U^{(n)}\to V
        \qquad\text{strongly in }L^3(Q_\theta),
\]
this proves
\[
        \dist_{\rm har}
        \bigl((U^{(n)},P^{(n)}),\calL^{\rm str}(Q_\theta)\bigr)
        \to0 .
\]
The proof is complete.
\end{proof}

\end{document}
